\chapter{1989年全国卷（文）}

\section{单选题}

\begin{problem}
如果 \(I = \{a, b, c, d, e\}\)，\(M = \{a, c, d\}\)，\(N = \{b, d, e\}\)，其中 \(I\) 是全集，那么 \(\overline{M} \cap \overline{N}\) 等于（\(\quad\)）

\choices
    {\(\varnothing\)}
    {\(\{d\}\)}
    {\(\{a, c\}\)}
    {\(\{b, e\}\)}
\end{problem}

\begin{answer}
A
\end{answer}

\begin{solution}
因为 \(I\) 是全集，所以 \(\overline{M}=I\mathbin{\backslash}M=\{b,e\}\)，\(\overline{N}=I\mathbin{\backslash}N=\{a,c\}\). 因此

\[\overline{M}\cap\overline{N}=\{b,e\}\cap\{a,c\}=\varnothing\]

故选 A.
\end{solution}

\begin{problem}
与函数 \(y = x\) 有相同图象的一个函数是（\(\quad\)）

\choices
    {\(y = \sqrt{x^2}\)}
    {\(y = \frac{x^2}{x}\)}
    {\(y = a^{\log_a x}\)，其中 \(a > 0\)，\(a \neq 1\)}
    {\(y = \log_a a^x\)，其中 \(a > 0\)，\(a \neq 1\)}
\end{problem}

\begin{answer}
D
\end{answer}

\begin{solution}
选项 A 中，\(\sqrt{x^2}=|x|\)，当 \(x<0\) 时不等于 \(x\). 选项 B 中，\(\dfrac{x^2}{x}=x\) 只在 \(x\neq0\) 时成立，定义域不包括 \(0\). 选项 C 中，\(a^{\log_a x}=x\) 的定义域为 \(x>0\)，也不是函数 \(y=x\) 的定义域.

对任意 \(x\in\R\)，都有 \(\log_a a^x=x\)，且 \(a>0\)，\(a\neq1\) 时该式有意义，因此选项 D 的函数与 \(y=x\) 的图象相同.
\end{solution}

\begin{problem}
如果圆锥的底面半径为 \(\sqrt{2}\)，高为 \(2\)，那么它的侧面积是（\(\quad\)）

\choices
    {\(4\sqrt{3}\pi\)}
    {\(2\sqrt{2}\pi\)}
    {\(2\sqrt{3}\pi\)}
    {\(4\sqrt{2}\pi\)}
\end{problem}

\begin{answer}
C
\end{answer}

\begin{solution}
圆锥的母线长为

\[l=\sqrt{\left(\sqrt{2}\right)^2+2^2}=\sqrt{6}\]

所以圆锥的侧面积为

\[\pi rl=\pi\cdot\sqrt{2}\cdot\sqrt{6}=2\sqrt{3}\pi\]

故选 C.
\end{solution}

\begin{problem}
已知 \(\{a_n\}\) 是等比数列，如果 \(a_1 + a_2 = 12\)，\(a_2 + a_3 = -6\)，且 \(S_n = a_1 + a_2 + \cdots + a_n\)，那么 \(\displaystyle\lim_{n \to \infty} S_n\) 的值等于（\(\quad\)）

\choices
    {\(8\)}
    {\(16\)}
    {\(32\)}
    {\(48\)}
\end{problem}

\begin{answer}
B
\end{answer}

\begin{solution}
设等比数列的公比为 \(q\). 由 \(a_2=a_1q\)、\(a_3=a_1q^2\)，有

\(a_1(1+q)=12\)，\(a_1q(1+q)=-6\).

因为 \(a_1(1+q)=12\neq0\)，两式相除得 \(q=-\dfrac12\). 于是 \(|q|<1\)，且

\[a_1=\frac{12}{1+q}=\frac{12}{1-\frac12}=24\]

因此无穷等比数列的和为

\[\lim_{n\to\infty}S_n=\frac{a_1}{1-q}=\frac{24}{1+\frac12}=16\]

故选 B.
\end{solution}

\begin{problem}
如果 \((1 - 2x)^7 = a_0 + a_1x + a_2x^2 + \cdots + a_7x^7\)，那么 \(a_1 + a_2 + \cdots + a_7\) 的值等于（\(\quad\)）

\choices
    {\(-2\)}
    {\(-1\)}
    {\(0\)}
    {\(2\)}
\end{problem}

\begin{answer}
A
\end{answer}

\begin{solution}
令 \(x=1\)，由展开式可得

\[a_0+a_1+a_2+\cdots+a_7=(1-2)^7=-1\]

令 \(x=0\)，可得 \(a_0=1\). 因而

\[a_1+a_2+\cdots+a_7=-1-a_0=-1-1=-2\]

故选 A.
\end{solution}

\begin{problem}
如果 \(\left|\cos \theta\right| = \frac{1}{5}\)，\(\frac{5\pi}{2} < \theta < 3\pi\)，那么 \(\sin \frac{\theta}{2}\) 的值等于（\(\quad\)）

\choices
    {\(-\frac{\sqrt{10}}{5}\)}
    {\(\frac{\sqrt{10}}{5}\)}
    {\(-\frac{\sqrt{15}}{5}\)}
    {\(\frac{\sqrt{15}}{5}\)}
\end{problem}

\begin{answer}
C
\end{answer}

\begin{solution}
因为 \(\frac{5\pi}{2}<\theta<3\pi\)，所以 \(\frac\pi2<\theta-2\pi<\pi\)，从而 \(\cos\theta<0\). 结合 \(\left|\cos\theta\right|=\frac15\)，得 \(\cos\theta=-\frac15\). 又因为 \(\frac{5\pi}{4}<\frac\theta2<\frac{3\pi}{2}\)，所以 \(\sin\frac\theta2<0\). 由半角公式，

\[\sin^2\frac\theta2=\frac{1-\cos\theta}{2}=\frac{1+\frac15}{2}=\frac35\]

因此 \(\sin\frac\theta2=-\frac{\sqrt{15}}5\)，故选 C.
\end{solution}

\begin{problem}
直线 \(2x + 3y - 6 = 0\) 关于点 \((1,-1)\) 对称的直线是（\(\quad\)）

\choices
    {\(3x - 2y + 2 = 0\)}
    {\(2x + 3y + 7 = 0\)}
    {\(3x - 2y - 12 = 0\)}
    {\(2x + 3y + 8 = 0\)}
\end{problem}

\begin{answer}
D
\end{answer}

\begin{solution}
设原直线上的点为 \((x,y)\)，它关于点 \((1,-1)\) 的对称点为 \((X,Y)\). 则

\(x=2-X\)，\(y=-2-Y\).

代入原直线方程，得

\[2(2-X)+3(-2-Y)-6=0\]

即 \(2X+3Y+8=0\). 因而所求直线为 \(2x+3y+8=0\)，故选 D.
\end{solution}

\begin{problem}
已知球的两个平行截面的面积分别为 \(5\pi\) 和 \(8\pi\)，它们位于球心的同一侧，且相距为 \(1\)，那么这个球的半径是（\(\quad\)）

\choices
    {\(4\)}
    {\(3\)}
    {\(2\)}
    {\(5\)}
\end{problem}

\begin{answer}
B
\end{answer}

\begin{solution}
设球的半径为 \(R\)，球心到两个截面的距离分别为 \(d_1\)，\(d_2\). 由于截面半径分别为 \(\sqrt5\) 和 \(2\sqrt2\)，有

\(d_1^2=R^2-5\)，\(d_2^2=R^2-8\).

面积为 \(5\pi\) 的截面半径较小，距离球心较远，因此 \(d_1-d_2=1\). 又

\[d_1^2-d_2^2=3\]

所以 \((d_1-d_2)(d_1+d_2)=3\)，从而 \(d_1+d_2=3\). 解得 \(d_1=2\)，\(d_2=1\)，进而

\(R^2=d_1^2+5=4+5=9\)，\(R=3\).

故选 B.
\end{solution}

\begin{problem}
由数字 \(1\)，\(2\)，\(3\)，\(4\)，\(5\) 组成没有重复数字的五位数，其中偶数共有（\(\quad\)）

\choices
    {\(60\text{ 个}\)}
    {\(48\text{ 个}\)}
    {\(36\text{ 个}\)}
    {\(24\text{ 个}\)}
\end{problem}

\begin{answer}
B
\end{answer}

\begin{solution}
五位数为偶数，当且仅当个位数字为偶数. 给定数字中个位可以取 \(2\) 或 \(4\)，有 \(2\) 种选法. 个位确定后，其余四个数字任意排列，有 \(4!\) 种排法. 因此偶数的个数为

\[2\times4!=2\times24=48\]

故选 B.
\end{solution}

\begin{problem}
如果双曲线 \(\frac{x^2}{64} - \frac{y^2}{36} = 1\) 上一点 \(P\) 到它的右焦点的距离是 \(8\)，那么点 \(P\) 到它的右准线的距离是（\(\quad\)）

\choices
    {\(10\)}
    {\(\frac{32\sqrt{7}}{7}\)}
    {\(2\sqrt{7}\)}
    {\(\frac{32}{5}\)}
\end{problem}

\begin{answer}
D
\end{answer}

\begin{solution}
双曲线的半实轴、半虚轴分别为 \(8\)、\(6\)，所以焦半距

\[c=\sqrt{8^2+6^2}=10\]

离心率为 \(e=\dfrac{c}{8}=\dfrac54\). 左支上任一点到右焦点 \((10,0)\) 的距离至少为 \(18\)，不可能等于 \(8\)，所以 \(P\) 在右支上. 由双曲线的定义，点到右焦点的距离与点到右准线的距离之比等于离心率. 设点 \(P\) 到右准线的距离为 \(d\)，则

\(\frac{8}{d}=\frac54\)，\(d=\frac{32}{5}\).

故选 D.
\end{solution}

\begin{problem}
如果 \(|x| \leqslant \frac{\pi}{4}\)，那么函数 \(f(x) = \cos^2 x + \sin x\) 的最小值是（\(\quad\)）

\choices
    {\(\frac{\sqrt{2} - 1}{2}\)}
    {\(\frac{1 + \sqrt{2}}{2}\)}
    {\(-1\)}
    {\(\frac{1 - \sqrt{2}}{2}\)}
\end{problem}

\begin{answer}
D
\end{answer}

\begin{solution}
令 \(t=\sin x\). 由 \(|x|\leqslant\dfrac\pi4\)，得

\[-\frac{\sqrt2}{2}\leqslant t\leqslant\frac{\sqrt2}{2}\]

又因为 \(\cos^2x=1-\sin^2x\)，所以

\[f(x)=1+t-t^2=\frac54-\left(t-\frac12\right)^2\]

这是关于 \(t\) 的开口向下的二次函数，故在上述闭区间上的最小值取在端点. 计算得

\[f\left(-\frac\pi4\right)=1-\frac12-\frac{\sqrt2}{2}=\frac{1-\sqrt2}{2}\]

\[f\left(\frac\pi4\right)=1-\frac12+\frac{\sqrt2}{2}=\frac{1+\sqrt2}{2}\]

所以最小值为 \(\dfrac{1-\sqrt2}{2}\)，故选 D.
\end{solution}

\begin{problem}
已知 \(f(x) = 8 + 2x - x^2\)，如果 \(g(x) = f(2 - x^2)\)，那么 \(g(x)\)（\(\quad\)）

\choices
    {在区间 \((-1,0)\) 上是减函数}
    {在区间 \((0,1)\) 上是减函数}
    {在区间 \((-2,0)\) 上是增函数}
    {在区间 \((0,2)\) 上是增函数}
\end{problem}

\begin{answer}
A
\end{answer}

\begin{solution}
由复合函数求导，

\[g'(x)=f'(2-x^2)(-2x)=\left(2-2(2-x^2)\right)(-2x)=4x(1-x^2)\]

在 \((-1,0)\) 上，\(x<0\) 且 \(1-x^2>0\)，所以 \(g'(x)<0\)，函数是减函数. 在 \((0,1)\) 上，\(g'(x)>0\)，函数是增函数.

在 \((-2,0)\) 和 \((0,2)\) 上，\(g'(x)\) 都会在 \(x=-1\) 或 \(x=1\) 处变号，不能分别保持单调递增. 因此只有选项 A 正确.
\end{solution}

\section{填空题}

\begin{problem}
给定三点 \(A(1,0)\)，\(B(-1,0)\)，\(C(1,2)\)，那么通过点 \(A\) 并且与直线 \(BC\) 垂直的直线方程是\fillinblank{}.
\end{problem}

\begin{answer}
\(x+y-1=0\)
\end{answer}

\begin{solution}
直线 \(BC\) 的斜率为

\[k_{BC}=\frac{2-0}{1-(-1)}=1\]

所求直线与 \(BC\) 垂直，所以其斜率为 \(-1\). 该直线经过 \(A(1,0)\)，故

\[y-0=-(x-1)\]

即 \(x+y-1=0\).
\end{solution}

\begin{problem}
不等式 \(|x^2 - 3x| > 4\) 的解集是\fillinblank{}.
\end{problem}

\begin{answer}
\(\left(-\infty,-1\right)\cup\left(4,+\infty\right)\)
\end{answer}

\begin{solution}
由绝对值不等式的定义，原不等式等价于

\[x^2-3x>4\quad\text{或}\quad x^2-3x<-4\]

第一种情况给出

\[x^2-3x-4>0\Longleftrightarrow (x-4)(x+1)>0\]

所以 \(x<-1\) 或 \(x>4\). 第二种情况给出

\[x^2-3x+4<0\]

但该二次式的判别式为 \(9-16=-7<0\)，且二次项系数为正，始终大于零，无解. 因此解集为 \(\left(-\infty,-1\right)\cup\left(4,+\infty\right)\).
\end{solution}

\begin{problem}
函数 \(y = \frac{\e^x - 1}{\e^x + 1}\) 的反函数的定义域是\fillinblank{}.
\end{problem}

\begin{answer}
\(\left(-1,1\right)\)
\end{answer}

\begin{solution}
令 \(t=\e^x\)，则 \(t>0\)，且原函数可写为 \(y=\dfrac{t-1}{t+1}\). 由于

\(y+1=\frac{2t}{t+1}>0\)，\(1-y=\frac{2}{t+1}>0\)

所以 \(-1<y<1\). 反过来，若 \(-1<y<1\)，则由 \(y=\dfrac{t-1}{t+1}\) 解得

\[t=\frac{1+y}{1-y}>0\]

确有相应的 \(x=\ln t\). 因此原函数的值域为 \(\left(-1,1\right)\)，也就是反函数的定义域.
\end{solution}

\begin{problem}
已知 \(A\) 和 \(B\) 是两个命题，如果 \(A\) 是 \(B\) 的充分条件，那么 \(B\) 是 \(A\) 的\fillinblank{}条件；\(\overline{A}\) 是 \(\overline{B}\) 的\fillinblank{}条件.
\end{problem}

\begin{answer}
必要；必要
\end{answer}

\begin{solution}
“\(A\) 是 \(B\) 的充分条件”表示 \(A\Rightarrow B\)，因此 \(B\) 是 \(A\) 的必要条件. 对这个命题取逆否命题，得到 \(\overline{B}\Rightarrow\overline{A}\)，所以 \(\overline{A}\) 是 \(\overline{B}\) 的必要条件.
\end{solution}

\begin{problem}
已知 \(0 < a < 1\)，\(0 < b < 1\)，\(a^{\log_b(x - 3)} < 1\)，那么 \(x\) 的取值范围是\fillinblank{}.
\end{problem}

\begin{answer}
\(3<x<4\)
\end{answer}

\begin{solution}
因为 \(0<a<1\)，指数函数 \(a^u\) 在 \(\R\) 上是减函数，且 \(a^0=1\)，所以

\[a^{\log_b(x-3)}<1\Longleftrightarrow \log_b(x-3)>0\]

又因为 \(0<b<1\)，对数函数 \(\log_b t\) 在 \(t>0\) 上是减函数，且 \(\log_b1=0\)，故

\[\log_b(x-3)>0\Longleftrightarrow 0<x-3<1\]

因此 \(x\) 的取值范围为 \(3<x<4\).
\end{solution}

\begin{problem}
如图，\(P\) 是二面角 \(\alpha - AB - \beta\) 棱 \(AB\) 上的一点，分别在 \(\alpha\)，\(\beta\) 上引射线 \(PM\)，\(PN\)，如果 \(\angle BPM = \angle BPN = 45^\circ\)，\(\angle MPN = 60^\circ\)，那么二面角 \(\alpha - AB - \beta\) 的大小是\fillinblank{}.

\bitmapfigure[max width=0.42\linewidth]{img/1989/national_liberal/q18_fig1.png}
\end{problem}

\begin{answer}
\(90^\circ\)
\end{answer}

\begin{solution}
取单位向量 \(\bs e\) 沿棱 \(PB\) 指向 \(B\)，取单位向量 \(\bs u\)、\(\bs v\) 分别垂直于 \(PB\) 且位于平面 \(\alpha\)、\(\beta\) 内. 设二面角的大小为 \(\varphi\)，则 \(\bs u\cdot\bs v=\cos\varphi\). 由两个夹角均为 \(45^\circ\)，可将两条射线的方向单位向量写成

\(\bs m=\frac{\sqrt2}{2}\bs e+\frac{\sqrt2}{2}\bs u\)，\(\bs n=\frac{\sqrt2}{2}\bs e+\frac{\sqrt2}{2}\bs v\).

于是

\[\cos60^\circ=\bs m\cdot\bs n=\frac12+\frac12\bs u\cdot\bs v=\frac12+\frac12\cos\varphi\]

所以 \(\cos\varphi=0\)，从而 \(\varphi=90^\circ\).
\end{solution}

\section{解答题}

\begin{problem}
设复数 \(z = \left(1 - \sqrt{3}\i\right)^5\)，求 \(z\) 的模和辐角的主值.
\end{problem}

\begin{solution}
复数 \(1-\sqrt3\i\) 的模为 \(2\)，其辐角的主值为 \(-\dfrac\pi3\)，所以

\[1-\sqrt3\i=2\left(\cos\left(-\frac\pi3\right)+\i\sin\left(-\frac\pi3\right)\right)\]

由棣莫弗定理，

\[z=2^5\left(\cos\left(-\frac{5\pi}{3}\right)+\i\sin\left(-\frac{5\pi}{3}\right)\right)=32\left(\cos\frac\pi3+\i\sin\frac\pi3\right)\]

因此 \(\lvert z\rvert=32\)，且 \(z\) 的辐角主值为 \(\dfrac\pi3\).
\end{solution}

\begin{problem}
证明：\(\tan \frac{3x}{2} - \tan \frac{x}{2} = \frac{2\sin x}{\cos x + \cos 2x}\).
\end{problem}

\begin{solution}
在等式两边有定义的范围内，利用正切函数作差公式，得

\[\tan\frac{3x}{2}-\tan\frac{x}{2}=\frac{\sin\left(\frac{3x}{2}-\frac{x}{2}\right)}{\cos\frac{3x}{2}\cos\frac{x}{2}}=\frac{\sin x}{\cos\frac{3x}{2}\cos\frac{x}{2}}\]

再由积化和差公式，

\[\cos\frac{3x}{2}\cos\frac{x}{2}=\frac{\cos2x+\cos x}{2}\]

代入上式即得

\[\tan\frac{3x}{2}-\tan\frac{x}{2}=\frac{2\sin x}{\cos x+\cos2x}\]

证毕.
\end{solution}

\begin{problem}
如图，在平行六面体 \(ABCD - A_1B_1C_1D_1\) 中，已知 \(AB = 5\)，\(AD = 4\)，\(AA_1 = 3\)，\(AB \perp AD\)，\(\angle A_1AB = \angle A_1AD = \frac{\pi}{3}\).

\begin{enumerate}
\item 求证：顶点 \(A_1\) 在底面 \(ABCD\) 的射影 \(O\) 在 \(\angle BAD\) 的平分线上；
\item 求这个平行六面体的体积.
\end{enumerate}

\bitmapfigure[max width=0.40\linewidth]{img/1989/national_liberal/q21_fig1.png}
\end{problem}

\begin{solution}
建立直角坐标系，使 \(A\) 为原点，\(AB\) 沿 \(x\) 轴正方向，\(AD\) 沿 \(y\) 轴正方向，底面 \(ABCD\) 为 \(z=0\). 设 \(A_1=(x,y,h)\)，其中 \(h>0\). 由两个夹角条件，

\(x=AA_1\cos\angle A_1AB=3\cos\frac\pi3=\frac32\)，\(y=AA_1\cos\angle A_1AD=3\cos\frac\pi3=\frac32\).

所以 \(A_1\) 在底面上的射影为 \(O\left(\dfrac32,\dfrac32,0\right)\). 又因为 \(AB\perp AD\)，底面内 \(\angle BAD\) 的平分线方向是 \((1,1)\)，故 \(O\) 在 \(\angle BAD\) 的平分线上.

由 \(AA_1=3\)，得

\[h=\sqrt{3^2-\left(\frac32\right)^2-\left(\frac32\right)^2}=\frac{3}{\sqrt2}\]

底面是矩形，面积为 \(AB\cdot AD=5\times4=20\). 因而平行六面体的体积为

\[V=20\cdot\frac{3}{\sqrt2}=30\sqrt2\]
\end{solution}

\begin{problem}
用数学归纳法证明：\((1 \cdot 2^2 - 2 \cdot 3^2) + (3 \cdot 4^2 - 4 \cdot 5^2) + \cdots + [(2n - 1)(2n)^2 - 2n(2n + 1)^2] = -n(n + 1)(4n + 3)\).
\end{problem}

\begin{solution}
记左端为 \(T_n\). 当 \(n=1\) 时，

\[T_1=1\cdot2^2-2\cdot3^2=4-18=-14=-1\cdot2\cdot7\]

等式成立.

假设当 \(n=k\) 时等式成立，即

\[T_k=-k(k+1)(4k+3)\]

则

\[\begin{gathered}T_{k+1}=T_k+(2k+1)(2k+2)^2-(2k+2)(2k+3)^2\\=-k(k+1)(4k+3)+(2k+2)\left((2k+1)(2k+2)-(2k+3)^2\right)\\=-k(k+1)(4k+3)-2(k+1)(6k+7)=-(k+1)(4k^2+15k+14)\\=-(k+1)(k+2)(4k+7).\end{gathered}\]

这正是 \(n=k+1\) 时右端 \(-(k+1)(k+2)\left(4(k+1)+3\right)\). 因此由数学归纳法，原等式对一切正整数 \(n\) 成立.
\end{solution}

\begin{problem}
已知 \(a > 0\)，\(a \neq 1\)，试求使方程 \(\log_a(x - ak) = \log_{a^2}(x^2 - a^2)\) 有解的 \(k\) 的取值范围.
\end{problem}

\begin{solution}
方程有意义必须满足

\(x-ak>0\)，\(x^2-a^2>0\).

又因为 \(\log_{a^2}u=\dfrac12\log_a u\)，原方程等价于

\[2\log_a(x-ak)=\log_a(x^2-a^2)\]

由于 \(a>0\) 且 \(a\neq1\)，对数函数 \(\log_a x\) 是一一对应的，所以

\[(x-ak)^2=x^2-a^2\]

整理得

\[-2akx+a^2k^2=-a^2\]

当 \(k=0\) 时上式变为 \(0=-a^2\)，不可能. 当 \(k\neq0\) 时，

\(x=\frac{a(k^2+1)}{2k}\)，\(x-ak=\frac{a(1-k^2)}{2k}\).

因为 \(a>0\)，定义域条件 \(x-ak>0\) 等价于

\[\frac{1-k^2}{k}>0\]

分类讨论：当 \(k>0\) 时，需 \(1-k^2>0\)，即 \(0<k<1\)；当 \(k<0\) 时，需 \(1-k^2<0\)，即 \(k<-1\). 对这些 \(k\)，由 \((x-ak)^2=x^2-a^2\) 且 \(x-ak>0\)，自动有 \(x^2-a^2>0\)，所以解确实满足全部定义域条件.

故 \(k\) 的取值范围为 \(\left(-\infty,-1\right)\cup\left(0,1\right)\).
\end{solution}

\begin{problem}
给定椭圆方程 \(\frac{x^2}{b^2} + \frac{y^2}{a^2} = 1\)（\(a > b > 0\)），求与这个椭圆有公共焦点的双曲线，使得以它们的交点为顶点的四边形面积最大，并求相应的四边形的顶点坐标.
\end{problem}

\begin{solution}
设椭圆焦半距为 \(c\)，则

\[c=\sqrt{a^2-b^2}\]

与椭圆有公共焦点的双曲线只能取纵轴为实轴的形式，设为 \(\frac{y^2}{u^2}-\frac{x^2}{v^2}=1\)，\(u>0\)，\(v>0\).

其焦半距满足 \(u^2+v^2=c^2\). 记两个曲线交点的横、纵坐标平方分别为 \(X=x^2\)、\(Y=y^2\)，则 \(\frac{X}{b^2}+\frac{Y}{a^2}=1\)，\(\frac{Y}{u^2}-\frac{X}{v^2}=1\).

解这个关于 \(X\)，\(Y\) 的方程组，得 \(X=\frac{b^2v^2(a^2-u^2)}{a^2v^2+b^2u^2}\)，\(Y=\frac{a^2u^2(b^2+v^2)}{a^2v^2+b^2u^2}\).

由 \(u^2+v^2=c^2=a^2-b^2\)，有 \(b^2+v^2=a^2-u^2\)，\(a^2v^2+b^2u^2=c^2(a^2-u^2)\).

因此

\[XY=\frac{a^2b^2u^2v^2}{c^4}\]

四个交点关于两坐标轴对称，组成的四边形为矩形，其面积为

\[S=4\sqrt{XY}=\frac{4abuv}{c^2}\]

由均值不等式，

\[uv\leqslant\frac{u^2+v^2}{2}=\frac{c^2}{2}\]

当且仅当 \(u=v=\dfrac{c}{\sqrt2}\) 时取等号. 所以面积最大时的双曲线为

\[\frac{y^2}{c^2/2}-\frac{x^2}{c^2/2}=1,\qquad\text{即}\qquad y^2-x^2=\frac{a^2-b^2}{2}\]

此时 \(X=\frac{b^2}{2}\)，\(Y=\frac{a^2}{2}\).

故四边形的四个顶点为 \(\left(\frac{b}{\sqrt2},\frac{a}{\sqrt2}\right)\)，\(\left(-\frac{b}{\sqrt2},\frac{a}{\sqrt2}\right)\)，\(\left(-\frac{b}{\sqrt2},-\frac{a}{\sqrt2}\right)\)，\(\left(\frac{b}{\sqrt2},-\frac{a}{\sqrt2}\right)\)

最大面积为 \(S_{\max}=2ab\).
\end{solution}
